The differences across the three environments help separate the role of fiscal scale from the role of the output distribution. Changing a constant surplus parameter rescales resources and optimal issuance, leaving the relative gain invariant. Changing ordinary-growth volatility or disaster incidence changes where the repayment schedule meets those resources. The larger percentage GDP-clause gain in the small-state illustration is therefore not generated mechanically by assigning it a lower $\alpha$. It reflects the different distribution of output around the common trigger, together with the change in the amount of straight debt that can profitably be issued. The stationary multiplier then maps those one-period factors into a further difference in recurring borrowing capacity.

The advanced-economy case supplies a complementary comparison. Its lower disaster probability reduces the contribution of a disaster-only index to insuring overall output stress. The GDP clause continues to respond to ordinary growth and logistics realizations. Its gain is consequently not a pure measure of disaster insurance: it also reflects the ability to alter repayments across non-disaster output states. This observation matters for the climate interpretation of the results. Climate exposure motivates the added tail process, but the general contract-design mechanism applies whenever the resources used to repay sovereign debt are stochastic.

The debt-limit table also helps interpret the GDP clause's optimum near $d=\alpha\bar g$ in the one-period problem. At that face, all states above the trigger can repay the full promise, while states below it face the lower promise. Moving face above the boundary can create default just above the trigger even when somewhat lower-output states still repay the reduced amount. The resulting survival set is no longer necessarily a single upper interval. The shape change is visible in Figure~\ref{fig:curves} and is established in Lemma~\ref{lem:threshold}; it is not a numerical artifact of interpolating a smooth default curve. A higher face therefore has to be assessed against the value of losing these above-trigger repayments.

Finally, the full-indexation bound identifies what remains outside the restricted designs. A step clause uses a small number of payment levels and cannot in general match repayment to resources in every realization. Full indexation does so under the maintained frictionless pricing rule. The bound is informative precisely because implementing that mapping may be difficult: the model places no charge on measurement, verification, incentive effects, or the institutional feasibility of pledging a large share of fiscal resources. Reporting the bound shows the distance between a simple contractual design and an unconstrained information benchmark; it does not recommend that the sovereign exhaust that benchmark.
